\begin{resumo}[Abstract]
\begin{otherlanguage*}{english}
This report describes a Solar System simulator for teaching physics, mathematics and computing, in which the instructor chooses how deep the explanation goes, from high school to graduate level. The project started from an undergraduate thesis that compared Copernicus' epicycles with the Fourier series of planetary orbits \cite{farias2016} and had three versions: a Python prototype that rebuilds the Kepler ellipse from a Fourier series; a two-dimensional version, in Python and Pygame, that integrates the mutual gravitation of the Sun and the eight planets from JPL Horizons initial conditions; and a three-dimensional version, in TypeScript and three.js, with thirteen moons, real imaging textures, a Fourier laboratory and a body-creation sandbox. The dynamics use Yoshida's fourth-order symplectic integrator with a fixed half-day step. Against JPL's reference ephemeris, the error stays below 0.5\,\% of the heliocentric distance on the five dates tested, from $-100$ to $+100$ years, and the energy oscillates around its initial value instead of drifting. The text covers the physical background (gravitation, the two-body problem, orbital elements, Kepler's equation), the mathematical background (complex Fourier series, epicycles, sampling, symplectic integrators), the implementation choices and the test-driven development process. On the thesis' question, we show, in agreement with analyses in the literature, that the complex Fourier series of a Keplerian orbit, to first order in the eccentricity, has the structure of Copernicus' construction: a fixed offset of $\tfrac{3}{2}ae$, a deferent of radius $a$ and an epicycle of radius $\tfrac{1}{2}ae$ turning at twice the frequency. The radii in the \emph{Commentariolus} for Mars, Jupiter and Saturn agree with these values to within 10\,\%. We did not compare this result with other models of the same order, such as Ptolemy's with equant, and the classroom uses suggested here have not been tested with students.

\vspace{\onelineskip}
\noindent
\textbf{Keywords}: epicycles. Fourier series. $N$-body problem. symplectic integrator. physics education.
\end{otherlanguage*}
\end{resumo}
